% Dedicated comparison report; not a replacement for the manuscript.
\documentclass{article}
\usepackage{iclr2027_conference}
\usepackage{amsmath,booktabs,tabularx,array,fontspec,float}
\setmainfont{texgyretermes-regular.otf}[Path=fonts/,BoldFont=texgyretermes-bold.otf,ItalicFont=texgyretermes-italic.otf,BoldItalicFont=texgyretermes-bolditalic.otf]
\title{BioCoLoop: Public-Harness Comparison}
\author{}
\date{}
\begin{document}
\maketitle
\noindent Public frameworks use their task-adapted native experimental controllers and the common biological fitting interface. In Panel B, AI-Scientist-v2 and AI-Researcher access laboratory 0, whereas BioCoLoop accesses ten laboratories. All Panel-B methods use Qwen2.5-7B-Instruct, the same twelve-design library and at most six candidate proposals. Scored cells pass independent recomputation from saved predictions; failure tokens record terminal controller outcomes rather than predictive scores.

\let\originaltable\table
\let\endoriginaltable\endtable
\renewenvironment{table}[1][]{\originaltable[H]}{\endoriginaltable}
\begin{table}[t]
\centering\scriptsize
\setlength{\tabcolsep}{1.55pt}
\renewcommand{\arraystretch}{0.82}
\begin{tabularx}{\linewidth}{@{}Xrrrrr@{}}
\toprule
\textbf{Model / method} & \multicolumn{1}{c}{DTI} & \multicolumn{1}{c}{Proteomics} & \multicolumn{3}{c}{Cell perturbation} \\
\cmidrule(lr){2-2}\cmidrule(lr){3-3}\cmidrule(l){4-6}
 & \shortstack{BindingDB\\AUROC $\uparrow$} & \shortstack{PTPC\\AP $\uparrow$} & \shortstack{VCC\\Top-1 $\uparrow$} & \shortstack{Norman\\Top-1 $\uparrow$} & \shortstack{Tahoe\\Top-1 $\uparrow$} \\
\midrule
\multicolumn{6}{l}{\textbf{Qwen2.5-7B-Instruct}} \\
Fixed model (1 lab) & 82.83 $\pm$ 1.22 & \underline{25.05} $\pm$ 4.56 & \underline{29.94} $\pm$ 0.00 & \textbf{25.56} $\pm$ 0.00 & \underline{19.40} $\pm$ 0.00 \\
Direct (1 lab) & \underline{89.89} $\pm$ 1.02 & 20.13 $\pm$ 0.07 & 21.09 $\pm$ 4.38 & 14.15 $\pm$ 0.00 & 15.42 $\pm$ 6.21 \\
AI-Scientist-v2 (1 lab) & 89.68 $\pm$ 0.80 & 20.18 $\pm$ 0.43 & 18.61 $\pm$ 2.78 & 14.12 $\pm$ 0.04 & 15.42 $\pm$ 6.21 \\
AI-Researcher (1 lab) & 86.09 $\pm$ 0.88 & 20.12 $\pm$ 0.05 & 19.70 $\pm$ 2.87,\,$F_{\mathrm{tool}}$\,(2/3) & 21.73 $\pm$ 6.63 & 17.41 $\pm$ 6.21 \\
\textbf{BioCoLoop} (10 labs) & \textbf{93.76} $\pm$ 1.11 & \textbf{36.99} $\pm$ 0.22 & \textbf{31.54} $\pm$ 0.08 & \underline{22.25} $\pm$ 0.45 & \textbf{22.39} $\pm$ 2.59 \\
\midrule
\multicolumn{6}{l}{\textbf{GPT-5.6 Luna (low reasoning)}} \\
Fixed model (1 lab) & 82.83 $\pm$ 1.22 & \underline{25.05} $\pm$ 4.56 & \underline{29.94} $\pm$ 0.00 & \textbf{25.56} $\pm$ 0.00 & \underline{19.40} $\pm$ 0.00 \\
Direct (1 lab) & \underline{89.89} $\pm$ 1.02 & 20.17 $\pm$ 0.42 & 18.61 $\pm$ 2.78 & 21.75 $\pm$ 6.59 & 15.92 $\pm$ 6.03 \\
AI-Scientist-v2 (1 lab) & 89.59 $\pm$ 0.95 & 23.37 $\pm$ 5.89 & 18.61 $\pm$ 2.78 & 14.12 $\pm$ 0.04 & 11.94 $\pm$ 2.11,\,$F_{\mathrm{svc}}$\,(2/3) \\
AI-Researcher (1 lab) & $F_{\mathrm{ctx}}$/$F_{\mathrm{svc}}$\,(0/3) & $F_{\mathrm{ctx}}$\,(0/3) & $F_{\mathrm{ctx}}$\,(0/3) & $F_{\mathrm{ctx}}$\,(0/3) & $F_{\mathrm{ctx}}$/$F_{\mathrm{svc}}$\,(0/3) \\
\textbf{BioCoLoop} (10 labs) & \textbf{93.67} $\pm$ 0.98 & \textbf{37.11} $\pm$ 0.01 & \textbf{31.50} $\pm$ 0.00 & \underline{22.69} $\pm$ 0.73 & \textbf{22.39} $\pm$ 2.59 \\
\bottomrule
\end{tabularx}
\caption{Main comparison by research model (mean $\pm$ sample SD, seeds 42--44). Public controllers use laboratory 0 and BioCoLoop uses ten; all rows share the design library, candidate cap, 100-round fits and held-out scorer. Scores are $\times100$; bold/underline mark the best/second-best complete result per block, and $(n/3)$ reports completed seeds without imputing failures.}
\label{tab:public-harness-comparison-three-seed}
\end{table}


\paragraph{Outcome.} BioCoLoop leads four of five endpoints: DTI (94.60\% AUROC), proteomics (37.12\% AP), VCC (31.63\% Top-1) and Tahoe (23.88\% Top-1). AI-Researcher leads Norman with 25.56\% Top-1, versus 22.74\% for BioCoLoop. AI-Scientist-v2's DTI tree-search terminates after three trained candidates when its native substage generator fails schema validation; the table reports $F_{\mathrm{pipe}}$ rather than an imputed score.
\end{document}
